% chapters/05/section-05.tex
\section{Program Gerak Lengkap}
\label{sec:siklus-gerak}

Bagian ini menyatukan seluruh unsur menjadi program yang menggerakkan objek
dengan kecepatan tetap, sekaligus mengukur kecepatan itu untuk membuktikan
bahwa nilainya tidak bergantung pada laju bingkai.

\subsection{Kode Lengkap}

\kodepython{python/ch05_gerak.py}{Gerak objek dengan vektor dan selang waktu}{lst:ch05-gerak}

\subsection{Membedah Bagian-Bagiannya}

\textbf{Fungsi \code{read_direction} memegang tahap masukan.} Fungsi itu
mengembalikan vektor arah berdasarkan tombol yang sedang ditahan. Perhatikan
bahwa vektor yang dikembalikan belum dinormalkan; penormalan dilakukan di
tempat lain agar fungsi ini tetap sederhana dan hanya menangani masukan.

\textbf{Fungsi \code{update_position} memegang tahap pembaruan.} Di dalamnya
terdapat tiga pekerjaan: menormalkan arah dan mengalikannya dengan kecepatan,
memperbarui posisi dengan perpindahan selama satu selang waktu, lalu membatasi
posisi agar objek tidak keluar dari jendela. Ketiganya berkaitan erat sehingga
diletakkan pada satu tempat.

\textbf{Penyelarasan \code{Rect} dilakukan di antara pembaruan dan
penayangan.} Dua baris berikut menjadi tempat terjadinya penyelarasan itu.

\begin{kodekeluaran}
previous_position = position.copy()
position, velocity = update_position(...)
step = position - previous_position
player_rect.center = (round(position.x), round(position.y))
\end{kodekeluaran}

Salinan posisi disimpan lebih dahulu agar perpindahan yang sesungguhnya dapat
dihitung sesudahnya. Perpindahan itulah yang diukur untuk membuktikan
kecepatan tetap, sebagaimana diuraikan pada Bagian~\ref{sec:gerak-selang-waktu}.

\textbf{Penayangan memperlihatkan kedua nilai sekaligus.}
Jendela menampilkan posisi pecahan dan posisi bulat secara berdampingan.

\begin{kodekeluaran}
Posisi (pecahan)   : ( 312.47,  240.00)
Rect (bulat)       : ( 312,     240)
\end{kodekeluaran}

Penampakan keduanya disengaja. Dengan mengamati kedua baris itu selama objek
bergerak, Anda dapat melihat bahwa nilai pecahan berubah secara halus
sedangkan nilai bulat melompat satu piksel setiap kali. Itulah wujud
penyelarasan yang dijelaskan pada Bagian~\ref{sec:selaras-rect}.

\subsection{Percobaan yang Disarankan}

\begin{table}[htbp]
\centering
\caption{Percobaan pada program gerak}
\label{tab:percobaan-ch05}
{\setlength{\tabcolsep}{4pt}\footnotesize%
\begin{tabularx}{\textwidth}{@{}L{5.6cm} Y@{}}
\toprule
\headrow \textbf{Yang diubah dan caranya} & \textbf{Yang sebaiknya Anda amati} \\
\midrule
Laju bingkai, melalui pilihan \code{--fps}. &
Kecepatan terukur tetap sama, sedangkan jumlah perpindahan setiap bingkai
berubah. Bandingkan nilai \code{delta_seconds} pada judul jendela. \\
\addlinespace
Kecepatan, melalui pilihan \code{--speed}. &
Kecepatan terukur berubah sebanding. Pada kecepatan sangat rendah, misalnya
5 piksel per detik, gerak objek menjadi tersendat karena pembulatan. \\
\addlinespace
Arah diagonal, dengan menekan dua tombol arah sekaligus. &
Kecepatan terukur tetap sama dengan gerak lurus. Inilah bukti bahwa penormalan
vektor arah bekerja. \\
\addlinespace
Baris penormalan, dengan menghapus pemanggilan \code{normalize()}. &
Kecepatan terukur pada arah diagonal naik menjadi sekitar 283 piksel per detik.
Kembalikan baris itu sesudah percobaan. \\
\addlinespace
Pembulatan, dengan mengganti \code{round()} menjadi \code{int()}. &
Selisih letak objek tidak terlihat pada pandangan sekilas, tetapi baris
\istilah{Rect} pada jendela memperlihatkan nilainya berbeda hampir satu
piksel. \\
\bottomrule
\end{tabularx}}
\end{table}

\begin{peringatanbox}
Setelah selesai bereksperimen, kembalikan program Anda ke keadaan semula.
Kebiasaan menyimpan keadaan percobaan pada repositori membuat riwayat
pekerjaan Anda tetap utuh, dan memudahkan Anda membandingkan perilaku program
sebelum dan sesudah percobaan. Inilah salah satu manfaat Git yang akan terasa
ketika pekerjaan Anda bertambah besar.
\end{peringatanbox}
